\originalchapter{Wilson 重整化群与临界现象}
\label{ch:wilson-rg}

传统重整化从发散图出发,Wilson 的出发点则是一个有限截止下的理论.我们通常只在某个能量区间知道哪些粒子和相互作用有效,没有理由先假定同一个作用量适用于任意短距离.若只测量低能过程,高能自由度虽然不会作为外部粒子出现,却会改变低能质量和顶角.重整化群就是把这种间接影响计算进新的作用量,并比较不同截止下的描述.

\originalsection{积分掉高动量自由度}

本章在 Euclidean 时空工作.设原截止为 $\Lambda_0$,单个实标量满足 $\phi\to-\phi$,作用量包含所有允许的局域项:
\begin{equation}
 S_{\Lambda_0}[\phi]=\int\dd^dx\left[
 \frac12(\partial\phi)^2+\frac12m_0^2\phi^2+
 \frac{\lambda_0}{4!}\phi^4+\frac{g_{6,0}}{6!}\phi^6+\cdots\right].
 \label{eq:wilson-bare-action}
\end{equation}
省略号包含更高次场和导数,不只是一个场幂级数.在固定 $\Lambda_0$ 下这些系数是有限输入,编码我们没有显式描述的高能物理.

选择 $\Lambda<\Lambda_0$,把 Fourier 分量分成支撑不相交的两部分:
\begin{equation}
 \phi=\phi_<+\phi_>,\quad
 \phi_<(k)=0\ (|k|>\Lambda),\quad
 \phi_>(k)=0\ (|k|<\Lambda\ \hbox{或}\ |k|>\Lambda_0).
\end{equation}
定义低能有效作用量为
\begin{equation}
 \e^{-S_\Lambda[\phi_<]}=
 \int\mathcal D\phi_>\,\e^{-S_{\Lambda_0}[\phi_<+\phi_>]}.
 \label{eq:wilson-effective-action}
\end{equation}
与场无关的归一化可独立放入真空常数.若源只含 $|k|<\Lambda$ 的动量分量,原路径积分就精确等于 $\int\mathcal D\phi_<\,\e^{-S_\Lambda+J\phi_<}$.这说明我们保留了什么:截止以下源所探测的关联函数,而不是把高动量场简单设为零.

积分产生的 $S_\Lambda$ 通常很复杂.平移不变性使每个动量顶角含 $\delta^{(d)}(\sum k_i)$;在外动量远小于已积分模式的尺度时,顶角可作导数展开,写成
\begin{equation}
 S_\Lambda=\int\dd^dx\left[\frac{Z(\Lambda)}2(\partial\phi_<)^2
                 +\sum_i g_i(\Lambda)\mathcal O_i(\phi_<)\right].
 \label{eq:wilson-local-expansion}
\end{equation}
把场重新归一化到动能系数一,会同时改变所有 $g_i$.完整作用量包含非局域动量依赖;局域展开是外动量相对截止小的展开,不能在外动量恰好等于截止时不加说明地截断到几个算符.光滑截止尤其适合这种展开.

先积分 $\Lambda<|k|<\Lambda_0$,再积分 $\Lambda'<|k|<\Lambda$,与一次积分整个区间相同.因此不同尺度的耦合形成有组合律的流.实际粗粒化丢失了高能信息,一般不能从任意低能截断作用量逆推出唯一裸作用;历史名称“群”在这一意义下常更接近半群.

\originalsection{工程维数与相关性的方向}

自由动能给 $[\phi]=(d-2)/2$.若 $[\mathcal O_i]=\Delta_i$,则 $[g_i]=d-\Delta_i\equiv\delta_i$.用截止定义无量纲耦合
\begin{equation}
 u_i(\Lambda)=g_i(\Lambda)\Lambda^{-\delta_i},\qquad
 \beta_i=\Lambda\frac{\dd u_i}{\dd\Lambda}.
 \label{eq:wilson-dimensionless}
\end{equation}
仅考虑工程缩放而忽略相互作用时,$g_i$ 不变,故
\begin{equation}
 \beta_i=-\delta_i u_i=(\Delta_i-d)u_i,
 \qquad u_i(\Lambda)=u_i(\Lambda_0)
                  (\Lambda/\Lambda_0)^{-\delta_i}.
 \label{eq:canonical-rg-flow}
\end{equation}
向红外是 $\Lambda$ 下降.$\Delta_i<d$ 时耦合增大,称相关;$\Delta_i>d$ 时耦合衰减,称无关;等号时称边缘,需要环修正决定方向.若改用向红外增加的参数 $t=\log(\Lambda_0/\Lambda)$,则 $\dd u_i/\dd t=-\beta_i$.后面的壳层计算采用 $t$;这一负号不能省略.

最直接的例子是自由高导数理论
\begin{equation}
 S=\frac12\int\frac{\dd^dk}{(2\pi)^d}
       \phi(k)\phi(-k)[m^2+k^2+a_4k^4+\cdots].
\end{equation}
低、高动量支撑不相交,积分掉高分量只产生常数,所以有量纲的 $m^2,a_4$ 都不改变.无量纲量却有 $r=m^2/\Lambda^2$、$v_4=a_4\Lambda^2$,分别满足
$\Lambda\dd r/\dd\Lambda=-2r$ 与 $\Lambda\dd v_4/\dd\Lambda=2v_4$.
这不是空洞的单位变化:在低动量核中,$m^2/k^2$ 确实越来越重要,而 $a_4k^4/k^2$ 越来越小.

若原截止处 $u_i(\Lambda_0)$ 是数量级一,在能量 $E\ll\Lambda_0$ 下,该项相对重要性按
$u_i(\Lambda_0)(E/\Lambda_0)^{\Delta_i-d}$ 变化.这解释了不可重整高维算符的压低,也指出相关质量项可能需要精细调节.这个工程论证以邻近 Gaussian 固定点、耦合尚未很强为前提;一般固定点要用实际缩放维数重新分类.

\originalsection{光滑截止下的精确流方程}

为了不把积分壳层变成每次都从头计算的任务,取光滑传播子
\begin{equation}
 G_\Lambda(k)=\frac{K(k^2/\Lambda^2)}{k^2},\qquad
 K(0)=1,\quad K(z)\to0\ (z\to\infty),
\end{equation}
这里 $K$ 取实值、非负且随 $z$ 单调不增,使 $G_\Lambda$ 与 $C=G_{\Lambda_0}-G_\Lambda$ 为半正定协方差;有限体积的常数零模先固定或排除,或者保留正的红外调节质量,最后仅在极限存在时移除调节器.并写 $S_\Lambda=\tfrac12\phi G_\Lambda^{-1}\phi+S_{I,\Lambda}$.质量项可放在 $S_I$ 中.为了使操作明确,可以先在有限体积、有限模式数下定义 Gaussian 测度,最后再讨论场论极限.

高模积分可看成协方差为 $C=G_{\Lambda_0}-G_\Lambda$ 的 Gaussian 卷积:
\begin{equation}
 \e^{-S_{I,\Lambda}[\phi]}
 =\int\dd\mu_C(\chi)\,\e^{-S_{I,\Lambda_0}[\phi+\chi]}.
 \label{eq:polchinski-gaussian-convolution}
\end{equation}
对一个有限维 Gaussian 卷积,协方差增加 $\delta C$ 的效果是
$\delta F=\tfrac12\delta C_{ab}\partial_a\partial_bF$.可把新增的小 Gaussian 增量 Taylor 展开来证明:其一次平均为零、二次平均为 $\delta C$,高阶属于更高微小量.由于 $\partial_\Lambda C=-\partial_\Lambda G_\Lambda$,有
\begin{equation}
 \Lambda\partial_\Lambda\e^{-S_I}
 =-\frac12\int\frac{\dd^dk}{(2\pi)^d}
     \Lambda\partial_\Lambda G_\Lambda(k)
     \frac{\delta^2\e^{-S_I}}{\delta\phi(k)\delta\phi(-k)}.
\end{equation}
利用 $\delta^2\e^{-S_I}=[(\delta S_I)^2-\delta^2S_I]\e^{-S_I}$,得到 Polchinski 形式的流方程
\begin{equation}
 \Lambda\partial_\Lambda S_I=
 \frac12\int\frac{\dd^dk}{(2\pi)^d}\,
 \Lambda\partial_\Lambda G_\Lambda(k)
 \left[\frac{\delta S_I}{\delta\phi(k)}
       \frac{\delta S_I}{\delta\phi(-k)}
       -\frac{\delta^2S_I}{\delta\phi(k)\delta\phi(-k)}\right].
 \label{eq:polchinski-flow}
\end{equation}
若使用未归一化的 Gaussian 测度,还会有由 $\tfrac12\Tr\log G_\Lambda$ 决定的真空常数流;它不改变这里的场依赖方程.

右边第一项把两个相互作用顶角用一条被积分的线连接,第二项在同一顶角上收缩两条腿.$\Lambda\partial_\Lambda G_\Lambda$ 主要支撑在 $|k|\sim\Lambda$,所以每步只处理一个薄壳.它只有看起来像树图和一圈图的结构,却是关于完整 $S_I$ 的方程;所有更高阶作用已经装在顶角中,没有作弱耦合展开.在有限调节器下,它是精确的函数方程.实际只保留有限种算符或有限导数时,才引入截断误差;“精确方程”不能自动保证某个有限截断的数值精确.

不同截止函数、场归一化与局域场重定义会给出不同的耦合坐标.例如 $\phi\to\phi+c\phi^3/\Lambda^2$ 同时改变导数四点项与六点项,但不能改变物理振幅.若 $\widetilde u_i=\widetilde u_i(u)$ 是可逆坐标变换,则
\begin{equation}
 \widetilde\beta_i=\sum_j\frac{\partial\widetilde u_i}{\partial u_j}\beta_j.
 \label{eq:beta-coordinate-vector}
\end{equation}
所以 beta 函数是耦合空间中的向量场,而单个坐标函数并不是方案无关的可观测量.由场重定义得到的冗余方向也不能误数成独立物理参数.

\originalsection{无关耦合为何仍会生成}

“无关”指初值的影响在红外衰减,不表示该耦合处处等于零.用四维的四点与六点耦合做一个局部模型.其流可具有
\begin{equation}
 \Lambda\frac{\dd u_4}{\dd\Lambda}=a u_6+\cdots,\qquad
 \Lambda\frac{\dd u_6}{\dd\Lambda}=2u_6-bu_4^2+\cdots,
 \label{eq:irrelevant-slaving-model}
\end{equation}
$a,b$ 的数值取决于算符归一化与截止方案.六点耦合的齐次解随 $\Lambda^2$ 衰减,但四点耦合不断给它源项.若 $u_4$ 变化缓慢,吸引的局部流形为
$u_6=(b/2)u_4^2+\order(u_4^3)$;把它代回第一式,得到
$\beta_4=(ab/2)u_4^2+\cdots$.
若该系数为正,$u_4$ 向红外缓慢下降,是边缘无关耦合.六点散射仍存在,只是其主导系数被四点耦合决定.这也说明传统只指定质量和四点耦合的可重整理论,不能被理解为低能有效作用中没有六点顶角.

有限个相关方向意味着要匹配有限个红外参数;无关方向的初值影响逐步被抹去.其余有效耦合则由流确定.在流保持于可控制区域时,这给出不同微观作用量产生相近低能行为的机制,也就是普适性.若流进入强耦合区,继续使用原场和有限截断可能失效,应寻找新的自由度,而不是仅凭无关初值衰减就宣称已证明整个红外理论.

\originalsection{固定点、缩放算符与临界面}

固定点 $u^*$ 满足所有 $\beta_i(u^*)=0$.Gaussian 固定点是无质量自由场,其他固定点可包含相互作用.固定点附近令 $\delta u_i=u_i-u_i^*$,线性化得
\begin{equation}
 \Lambda\frac{\dd\delta u_i}{\dd\Lambda}
   =\sum_jT_{ij}\delta u_j+\order(\delta u^2),\qquad
 T_{ij}=\left.\frac{\partial\beta_i}{\partial u_j}\right|_{u^*}.
\end{equation}
若 $e^{(a)}$ 是 $T$ 的左特征向量,$e^{(a)}T=t_a e^{(a)}$,则缩放变量
$U_a=\sum_ie_i^{(a)}\delta u_i$ 满足
\begin{equation}
 U_a(\Lambda)=U_a(\Lambda_0)(\Lambda/\Lambda_0)^{t_a}
 \quad\hbox{至线性阶}.
\end{equation}
所以 $y_a=-t_a>0$ 的方向相关,$y_a<0$ 的方向无关,$y_a=0$ 要检查非线性项.与 $U_a$ 共轭的缩放算符维数是 $d-y_a=d+t_a$.一般的原始算符会混合,不再逐个按工程维数分类.

在双曲固定点附近,即线性谱没有零实部方向时,稳定流形由无关方向的切空间决定;这就是局部临界面.要流向该固定点,必须把相关方向调到这个面上.边缘方向需要额外分析,临界面的曲率也来自非线性流,不能把线性切平面当作完整曲面.这里使用通常的稳定流形局部结论,不对无限维函数空间给出存在性证明.

\begin{center}
\begin{tikzpicture}[>=Stealth,scale=.95]
 \draw[->](-2.2,0)--(2.2,0) node[right] {$U_{\rm rel}$};
 \draw[->](0,-1.7)--(0,1.7) node[above] {$U_{\rm irr}$};
 \draw[very thick,structurecolor](0,-1.45)--(0,1.45);
 \draw[->,structurecolor](0,1.3)--(0,.5);
 \draw[->,structurecolor](0,-1.3)--(0,-.5);
 \draw[->] (.4,0)--(1.6,0);\draw[->](-.4,0)--(-1.6,0);
 \fill(0,0)circle(2pt);\node[anchor=west] at (.18,.28){$u^*$};
 \node[anchor=west,structurecolor] at (.35,1.25){临界面在此处的切线};
 \node at (0,-2.1){箭头表示向红外演化};
\end{tikzpicture}
\end{center}

固定点表示尺度不变性,但仅凭 $\beta=0$ 不应在所有维数和所有理论中不加条件地推出完整共形不变性.若相关扰动产生质量,红外也可能成为有质量理论;若耦合持续增大,则可能产生束缚态与新的低能场.后面 QCD 的一圈流只能辨认强耦合尺度,不能完成后一种情形的非微扰解.

\originalsection{四维附近的壳层计算}

为把概念落实为可复算的流,考虑 $d=4-\epsilon$,$0<\epsilon\ll1$ 的标量作用量
\begin{equation}
 S=\int\dd^dx\left[\frac12(\partial\phi)^2+
              \frac12m^2\phi^2+\frac\lambda{4!}\phi^4\right].
\end{equation}
这里只做弱耦合、低阶导数的 Wilson 截断,使用球形薄壳
$\Lambda\e^{-\dd t}<|k|<\Lambda$.
设
\begin{equation}
 K_d=\frac{\Omega_{d-1}}{(2\pi)^d},\quad K_4=\frac1{8\pi^2},
 \qquad r=\frac{m^2}{\Lambda^2},\quad
 g=K_d\lambda\Lambda^{-\epsilon}.
 \label{eq:phi4-shell-normalization}
\end{equation}
$t$ 向红外增加.下面的 $g$ 与前面的有量纲 $\lambda$ 不同,且与下一节采用的 $\lambda/(16\pi^2)$ 也相差因子二.

写 $\phi=\phi_<+\chi$,壳中 $\chi$ 的协方差是 $1/(k^2+m^2)$.累积量展开给
\begin{equation}
 \delta S=\langle S_I\rangle_>
   -\frac12\left(\langle S_I^2\rangle_>
                    -\langle S_I\rangle_>^2\right)+\cdots.
\end{equation}
一次项中 $6\phi_<^2\chi^2$ 产生 $\lambda\phi_<^2\langle\chi^2\rangle/4$,所以
\begin{equation}
 \delta m^2=\frac\lambda2 I_1,\qquad
 I_1=\int_{\rm shell}\frac{\dd^dk}{(2\pi)^d}\frac1{k^2+m^2}
    =\frac{K_d\Lambda^{d-2}}{1+r}\dd t+\order(\dd t^2).
\end{equation}
四点修正来自两个 $\phi_<^2\chi^2$ 项.连通收缩
$\langle\chi(x)^2\chi(y)^2\rangle_c=2G_>(x-y)^2$ 给系数
$-\lambda^2\phi_<^2(x)\phi_<^2(y)/16$.在低外动量下令两点的慢场相同,乘回 $4!=24$,便有
\begin{equation}
 \delta\lambda=-\frac32\lambda^2I_2,\qquad
 I_2=\int_{\rm shell}\frac{\dd^dk}{(2\pi)^d}\frac1{(k^2+m^2)^2}
 =\frac{K_d\Lambda^{d-4}}{(1+r)^2}\dd t+\order(\dd t^2).
\end{equation}
这既解释了负号来自第二累积量,也解释了三种四点通道给出的 $3/2$.一圈蝌蚪不依赖外动量,所以动能在这一阶不改变;场反常维数从两圈开始.

积分壳后将截止与动能恢复到相同单位,工程缩放贡献 $2r$ 与 $\epsilon g$,得到
\begin{align}
 \frac{\dd r}{\dd t}&=2r+\frac{g}{2(1+r)}+\order(g^2),\notag\\
 \frac{\dd g}{\dd t}&=\epsilon g-
             \frac{3g^2}{2(1+r)^2}+\order(g^3),\qquad
 \eta=\order(g^2).
 \label{eq:phi4-shell-flow}
\end{align}
壳层质量包含截止依赖的加性移位,所以临界点不在 $r=0$.求弱耦合固定点得到
\begin{equation}
 g^*=\frac{2\epsilon}{3}+\order(\epsilon^2),\qquad
 r^*=-\frac\epsilon6+\order(\epsilon^2).
 \label{eq:wilson-fisher-fixed-point}
\end{equation}
$r^*$ 随截止方案变化;固定点附近的缩放本征值才是比较物理的量.线性化 \eqref{eq:phi4-shell-flow},质量方向的本征值为
\begin{equation}
 y_t=2-\frac{g^*}{2}+\order(\epsilon^2)
      =2-\frac\epsilon3+\order(\epsilon^2),
\end{equation}
四点方向的本征值为 $-\epsilon+\order(\epsilon^2)$.虽然质量与四点耦合存在混合,反向混合从 $g^{*2}$ 开始,因此不会改变这里的一阶本征值.此固定点称为 Wilson--Fisher 固定点.

\originalsection{从流推出临界指数}

将系统离临界面的小距离记为温度型缩放变量 $\tau$.线性流给 $\tau(t)=\e^{y_tt}\tau(0)$.当 $|\tau(t_*)|$ 到数量级一时,固定点附近的尺度不变描述停止,相关长度按缩放恢复为 $\xi\sim\Lambda_0^{-1}\e^{t_*}$.因此
\begin{equation}
 \xi\propto|\tau|^{-\nu},\qquad
 \nu=\frac1{y_t}=\frac12+\frac\epsilon{12}+\order(\epsilon^2).
 \label{eq:phi4-critical-nu}
\end{equation}
固定点两点函数在 $\Lambda_0^{-1}\ll|x|\ll\xi$ 中按
$G(x)\sim|x|^{-(d-2+\eta)}$ 衰减.此处一圈 $\eta=0+\order(\epsilon^2)$,并非断言相互作用理论的场维数永远不变.

外场 $h$ 耦合于 $\phi$,故其本征值是 $y_h=(d+2-\eta)/2$.若临界奇异自由能满足通常的缩放形式
\begin{equation}
 f_s(\tau,h)=b^{-d}f_s(b^{y_t}\tau,b^{y_h}h),
 \label{eq:critical-free-energy-scaling}
\end{equation}
选 $b=|\tau|^{-1/y_t}$ 后求 $h$ 导数,得到磁化指数与易感率指数
\begin{equation}
 \beta_{\rm mag}=\frac{d-y_h}{y_t}
       =\frac12-\frac\epsilon6+\order(\epsilon^2),\qquad
 \gamma_{\rm sus}=\frac{2y_h-d}{y_t}
       =1+\frac\epsilon6+\order(\epsilon^2).
\end{equation}
这里 $\beta_{\rm mag}$ 是临界指数,不能与 RG 的 beta 函数混淆.再取 $\tau=0$、$b=|h|^{-1/y_h}$,磁化按 $h^{1/\delta_{\rm mag}}$ 变化,故
\begin{equation}
 \delta_{\rm mag}=\frac{y_h}{d-y_h}=3+\epsilon+\order(\epsilon^2),\qquad
 \alpha_{\rm heat}=2-d\nu=\frac\epsilon6+\order(\epsilon^2).
\end{equation}
$\gamma_{\rm sus}=(2-\eta)\nu$ 也可以直接由积分两点函数到 $|x|\sim\xi$ 得到.自由能的缩放假定没有破坏该形式的危险无关变量;在上临界维数以上,不应直接沿用这些 hyperscaling 关系.$\epsilon$ 展开给的是小 $\epsilon$ 的渐近微扰结论,把 $\epsilon=1$ 代入只能作为三维的低阶估计,不能宣称误差已经受控.

\originalsection{两个实标量的耦合流}

多耦合例子能检验“原来的算符不一定是缩放方向”这一点.考虑独立的 $\mathbb Z_2\times\mathbb Z_2$ 对称性,Euclidean 势为
\begin{equation}
 V=\frac12m_1^2\phi_1^2+\frac12m_2^2\phi_2^2+
 \frac1{4!}\left(\lambda_1\phi_1^4+\lambda_2\phi_2^4
                    +6\lambda_{12}\phi_1^2\phi_2^2\right).
 \label{eq:two-scalar-rg-potential}
\end{equation}
如此定义混合顶角的大小就是 $\lambda_{12}$.稳定的四次势至少要求 $\lambda_1,\lambda_2\geq0$ 且
$\lambda_{12}\geq-\sqrt{\lambda_1\lambda_2}/3$;否则还需要高次项才能使作用量有下界.

将有量纲耦合写为全对称张量 $\lambda_{abcd}\phi_a\phi_b\phi_c\phi_d/4!$,有
$\lambda_{1111}=\lambda_1$、$\lambda_{2222}=\lambda_2$、任意两个 $1$ 两个 $2$ 的分量为 $\lambda_{12}$.定义无量纲张量 $\bar\lambda_{abcd}=\lambda_{abcd}\Lambda^{-\epsilon}$.一圈四点图的三个通道给
\begin{equation}
 \beta_{abcd}\equiv\Lambda\frac{\dd\bar\lambda_{abcd}}{\dd\Lambda}
 =-\epsilon\bar\lambda_{abcd}+\frac1{16\pi^2}
 \sum_{m,n}\left(\bar\lambda_{abmn}\bar\lambda_{mncd}
 +\bar\lambda_{acmn}\bar\lambda_{mnbd}
 +\bar\lambda_{admn}\bar\lambda_{mnbc}\right).
 \label{eq:two-scalar-tensor-beta}
\end{equation}
此式的系数由前述每通道泡图恢复:$K_4/2=1/(16\pi^2)$.对 $1111$ 分量,每通道的求和是 $\lambda_1^2+\lambda_{12}^2$;对 $1122$,第一个通道是 $\lambda_{12}(\lambda_1+\lambda_2)$,其余两个各含两种内部物种排序,共给 $4\lambda_{12}^2$.

为使系数简洁,本节改定义
$g_i=\lambda_i\Lambda^{-\epsilon}/(16\pi^2)$、$w=\lambda_{12}\Lambda^{-\epsilon}/(16\pi^2)$.近无质量区、向红外的流为
\begin{align}
 \frac{\dd g_1}{\dd t}&=\epsilon g_1-3(g_1^2+w^2),\notag\\
 \frac{\dd g_2}{\dd t}&=\epsilon g_2-3(g_2^2+w^2),\notag\\
 \frac{\dd w}{\dd t}&=\epsilon w-w(g_1+g_2)-4w^2,
 \label{eq:two-scalar-coupling-flow}
\end{align}
至一圈.离开临界质量区时需要恢复质量分母,不能把此式用于任意大质量.

解耦固定点为 $g_1=g_2=\epsilon/3,w=0$.在它附近,混合扰动满足
$\dd w/\dd t=(\epsilon/3)w$,所以是相关的;两个单独的 Ising 固定点并不对一切允许混合稳定.另一方面,$O(2)$ 对称性对应
$\lambda_1=\lambda_2=u$、$\lambda_{12}=u/3$.在这个子空间中,固定点是
\begin{equation}
 g_1^*=g_2^*=\frac{3\epsilon}{10},\qquad
 w^*=\frac\epsilon{10}.
 \label{eq:two-scalar-o2-fixed}
\end{equation}
代回三条方程可逐项验证.

其四次耦合的稳定矩阵为
\begin{equation}
 \frac1\epsilon\left.\frac{\partial(\dot g_1,\dot g_2,\dot w)}
        {\partial(g_1,g_2,w)}\right|_*
 =\begin{pmatrix}-4/5&0&-3/5\\0&-4/5&-3/5\\-1/10&-1/10&-2/5\end{pmatrix}.
\end{equation}
反对称方向 $g_1-g_2$ 的本征值是 $-4\epsilon/5$;在对称子空间中剩下的本征值为 $-\epsilon$ 与 $-\epsilon/5$.所以这三个四次扰动在一圈、小 $\epsilon$ 下都向红外衰减.它们的线性组合才是缩放变量,不能单凭 $g_1,g_2,w$ 的工程维数判断.

质量方向仍然相关.壳层蝌蚪给
\begin{equation}
 \dot r_1=2r_1+\frac{g_1}{1+r_1}+\frac w{1+r_2}+\cdots,
 \qquad
 \dot r_2=2r_2+\frac{g_2}{1+r_2}+\frac w{1+r_1}+\cdots.
\end{equation}
在 \eqref{eq:two-scalar-o2-fixed} 附近,质量和与质量差的本征值分别为
$y_+=2-2\epsilon/5$、$y_-=2-\epsilon/5$.
若限制 $O(2)$ 对称性,只需调节质量和,相应
$\nu_{O(2)}=1/2+\epsilon/10+\order(\epsilon^2)$;在仅有 $\mathbb Z_2\times\mathbb Z_2$ 的多临界问题中,两个质量方向都需要调节.

这些计算示范了多个固定点、算符混合和不同临界面,不是完整相图分类.流若离开稳定弱耦合区域,可能提示所选截断不足,不能仅凭一圈“跑向无穷”的曲线就证明发生一阶相变.本节的壳层临界指数与两场计算是对原课固定点框架的教学展开;原课后半还会用传统重整化方案得到相同的尺度信息.
